\documentclass{wlemph}
\usepackage[utf8]{inputenc}
\usepackage[T1]{fontenc}
\usepackage{geometry}
\usepackage{graphicx}
\usepackage{float}
\usepackage{hyperref}
\usepackage[english,indonesian]{babel}

\title{Optimasi Pemrosesan Template LaTeX Menggunakan Pendekatan Berbasis AI}
\author{Academic Support Agent}
\date{\today}

\begin{document}
\maketitle

\begin{abstract}

\end{abstract}

\chapter{PENDAHULUAN}


\chapter{TINJAUAN PUSTAKA}


% Original content (commented out):
% \documentclass[fleqn]{wlemph}

\title{EMPH Clinical Brief}
\author[1]{First Author Name}
\author[2]{Second Author Name}
\author[3]{Third Author Name}
\affil[1]{First Author affiliation}
\affil[2]{Second Author affiliation}
\affil[3]{Third Author affiliation}

\begin{document}

\flushbottom
\maketitle
\thispagestyle{empty}
\noindent\begin{minipage}[t]{0.31333\textwidth}\setlength{\parindent}{15pt}

\section*{Name of pathology/condition}
\fontsize{9.5}{12}\selectfont

\textbf{Instructions for Authors}

You must use the titles for each column provided above. The first column should describe the pathology or condition from a clinical perspective.  

Fill this and the other two columns with appropriate text reflecting the column headings. Please use the type and size provided.

Use accessible language that can be easily understood by non-specialists. Avoid jargon. 

The total word limit for the contents of thee three columns is 600 words. The text must not exceed this limit. This includes headings, text, references and figures (30 words must be deducted for each figure), but not the article title and author details (above).

\end{minipage}%
\kern.03\textwidth%
\begin{minipage}[t]{0.31333\textwidth}\setlength{\parindent}{15pt}

\section*{Evolutionary perspectives}
\fontsize{9.5}{12}\selectfont

The second column should give an evolutionary perspective on this pathology or condition.  Provide alternative hypotheses where possible, and avoid advocacy.  
 
Contributions may have clear, explanatory figures or diagrams, as appropriate, in any or all of the text boxes, up to a maximum of 3. Please note that figures should not occupy more than a third of a full-length column, and 30 words must be deducted from the total word limit for each figure.

To include a figure, please upload your file using the project menu in the Overleaf editor, and edit the example figure box at the end of this template.  You may also need to submit a high-resolution version of the figure separately. Refer to our \href{http://www.oxfordjournals.org/for_authors/figures.html}{figure submission instructions} for further information about formatting and permissions for reproduced figures.

\end{minipage}%
\kern.03\textwidth%
\begin{minipage}[t]{0.31333\textwidth}\setlength{\parindent}{15pt}

\section*{Future implications}
\fontsize{9.5}{12}\selectfont

The final column should discuss the implications of an evolutionary perspective, and how this might improve clinical practice.  The author(s) should specify evidence base, levels of confidence and if, for example, there have been any clinical trials testing evolutionary hypotheses.

LaTeX formats citations and references automatically using the bibliography records in your .bib file (see example below). To add your own references, open the sample.bib file via the project menu and follow the examples provided. Use the cite command for an inline citation, e.g. articles \cite{Figueredo:2009dg} and \cite{mercer1985effect}, books \cite{Sutton1999Reinforcement} and websites \cite{Xithalis2011Hermes}.

Up to eight references can be added below the text in any of the columns depending on text length, formatting, and space.

When you are ready to submit your manuscript, please click the ``Submit to EMPH'' button on the topbar of the Overleaf editor, and follow the instructions provided. We hope you find Overleaf useful for your EMPH submission, and please \href{https://www.overleaf.com/contact}{let us know} if you have any feedback.

\bigskip
\bibliography{sample}

\end{minipage}%

% * <john.hammersley@gmail.com> 2015-02-25T16:30:56.499Z:
%
%  If you do not wish to include a figure, please delete the text block below, including the \begin{textblock} and \end{textblock} commands.
%

\begin{textblock}{2}(0,5) % Controls for {sizing} and (positioning) figure 
% * <john.hammersley@gmail.com> 2015-01-30T16:50:24.055Z:
%
%  The textblock command above controls the size and position of the figure. The first number controls the width, e.g. {2} specifies a figure to span 2 columns. The second pair of numbers (x,y) specifies the position. Please select 0,1 or 2 for x to position the figure with it's left hand edge in column 1, 2 or 3 respectively. Please adjust y as required to specify the correct vertical position.
%
\begin{figure}
\captionsetup{type=figure}
\includegraphics[width=\textwidth]{Fig.png}
\caption{\label{fig:label}To upload an image please use the `Add files' option in the project menu. Once uploaded, click the image preview in the editor and edit the \textit{includegraphics} line with the name of your file.}
\end{figure}
\end{textblock}

\end{document}

\end{document}